\section{Conclusion}

We studied dense full materialization for tensor networks with open legs. Unlike conventional approximate contraction, which typically targets scalar outputs or local queries, our setting requires the final result to be emitted explicitly as a dense tensor. Under this requirement, the output representation cannot be changed, so the natural place for approximation is the internal contraction process rather than the output side. Based on this view, we introduced \method{}.

\method{} organizes full materialization as a unified pipeline: it partitions the output index space, builds a partition tree on the interaction graph, recursively propagates separator states while keeping the output boundary explicit, and suppresses interface growth through adaptive low-rank compression and implicit sketch-based merging. The method therefore approximates not the final output object, but the internal separator states that must be traversed in order to produce it.

Experiments on Ring-8, Tree-8, Grid $3\times 3$, and Random-8 show that \method{} delivers a clear and usable accuracy-time trade-off, achieving substantial speedups over exact contraction while keeping relative error small. Stage-wise decomposition shows that the gains come almost entirely from the contraction stage. Mechanism analysis identifies implicit sketch-based merging as the most important and stable efficiency source, with output blocking also contributing on harder networks. Strength scans further show that the approximation level provides an effective knob for balancing accuracy and runtime.

More broadly, the results support a simple message: when the dense output is part of the task definition, the right target for approximation is not the output tensor itself, but the internal separator states propagated during contraction. Along this direction, partition-tree recursion, separator-state representation, and sketch-driven compression can be combined into a practical framework for open-leg tensor network contraction with full materialization.
